\scanpage{26}
Rédaction n\textsuperscript{o} 397

COMPL\handwritten{\'}{E}MENTS AU CHAPITRE VII D'ALG\handwritten{\`}{E}BRE

DUALIT\handwritten{\'}{E} DES MODULES DE LONGUEUR FINIE SUR UN ANNEAU PRINCIPAL

\begin{center}-:-:-:-\end{center}

1. \uline{Modules injectifs}.

DEFINITION 1.- \uline{Soit} $A$ \uline{un anneau}, \uline{et soit} $I$ \uline{un} $A$-\uline{module}. \uline{On dit que} $I$ \uline{est injectif si}, \uline{pour toute suite exacte} $E'\longrightarrow E\longrightarrow E''$ \uline{d'applications} $A$-\uline{linéaires}, \uline{la suite} :
\[
\operatorname{Hom}(E'',I)\longrightarrow\operatorname{Hom}(E,I)\longrightarrow\operatorname{Hom}(E',I)
\]
\uline{est exacte}.

PROPOSITION 1.- \uline{Soit} $I$ \uline{un} $A$-\uline{module}. \uline{Les conditions suivantes sont équivalentes} :

(i) \quad $I$ \uline{est injectif}.

(ii) \quad \uline{Pour tout homomorphisme injectif} $F\longrightarrow G$, \uline{l'homomorphisme correspondant} $\operatorname{Hom}(G,I)\longrightarrow\operatorname{Hom}(F,I)$ \uline{est surjectif}.

Posons, pour simplifier l'écriture, (et pour tricher un peu !) :
\[
D(F)=\operatorname{Hom}(F,I).
\]
La condition (ii) signifie que, si $0\longrightarrow F\longrightarrow G$ est une suite exacte, \deleted{\unsure{L|l}}\handwritten{l}a suite :
\[
D(G)\longrightarrow D(F)\longrightarrow 0
\]
est exacte. On voit bien qu'elle est entraînée par (i).

Inversement, supposons (ii) vérifiée, et soit $E'\longrightarrow E\longrightarrow E''$ une suite exacte quelconque. Considérons la décomposition canonique de cette suite, (cf. Alg. \unsure{'}II, 3e édition, p. 21) :
\diagram{\begin{tikzcd}[ampersand replacement=\&,column sep=1.4em,row sep=.4em]
\& E'\arrow[r] \& S\arrow[r] \& 0 \& {}\\
0\arrow[r] \& S\arrow[r] \& E\arrow[r] \& T\arrow[r] \& 0\\
\& 0\arrow[r] \& T\arrow[r] \& E''\text{.} \& {}
\end{tikzcd}}
% archive sheet number: 25

